\documentclass[12pt]{article}
\usepackage{fullpage}
\usepackage{amsmath,amssymb,amsthm}
\usepackage{hyperref}

\newtheorem{theorem}{Theorem}
\newtheorem{lemma}{Lemma}
\newtheorem{proposition}{Proposition}
\newtheorem{corollary}{Corollary}
\theoremstyle{definition}
\newtheorem{definition}{Definition}
\newtheorem{remark}{Remark}

\DeclareMathOperator{\sgn}{sgn}

\begin{document}

\begin{center}
\textbf{\Large The Hirzebruch signature theorem and the functional equation of the
$L$--genus}
\end{center}

\noindent\textbf{Overall status (read first).}
\emph{Part (b)} is proved completely and elementarily below
(\S\ref{sec:partb}): the functional equation of $\coth$ (Theorem~\ref{thm:funeq}), the
characterization of $z\coth z$ as the unique normalised genus equal to the signature on
all $\mathbb{CP}^{2m}$ (Theorem~\ref{thm:char}), and Novikov additivity
(Theorem~\ref{thm:novikov}) are all established from first principles; the
fibre-bundle incarnation (Chern--Hirzebruch--Serre, Theorem~\ref{thm:CHS}) is stated with
verified hypotheses and cited, with its elementary product case proved
(Proposition~\ref{prop:prodmult}). \emph{Part (a)} is reduced, by entirely elementary
arguments (Poincar\'e--Lefschetz duality, linear algebra, explicit Pontryagin-class
computations), to the \textbf{single} classical input
$\dim_{\mathbb Q}(\Omega^{SO}_{4k}\otimes\mathbb Q)\le p(k)$ (Thom's rational cobordism
rank, $(\sharp)$ below). We prove everything else --- cobordism invariance of the
signature, vanishing of Pontryagin numbers on boundaries, the spanning family with
closed-form values, and the deduction $\tau=L[\,\cdot\,]$ --- unconditionally and
elementarily, and we isolate $(\sharp)$ as the precise irreducible gap. As discussed in
Remark~\ref{rem:honest}, $(\sharp)$ is essentially the content the problem forbids
(cobordism theory) and no genuinely elementary proof of it is known; we therefore
present part (a) as the strongest complete reduction rather than an unconditional
elementary proof, and do not overclaim.

\bigskip

\begin{center}
\textbf{\large Part (a): a rationally generating family of $4k$--manifolds, by direct
linear algebra}\\[2pt]
\textit{(self-contained resolution of the ``generation gap'')}
\end{center}


\bigskip

\noindent
\textbf{Scope.}
This section supplies, by elementary means, the linear-algebra step that the standard
proof of the Hirzebruch signature theorem usually imports from cobordism theory
(Thom's computation $\Omega^{SO}_\ast\otimes\mathbb{Q}=
\mathbb{Q}[[\mathbb{CP}^2],[\mathbb{CP}^4],\dots]$).  We never use the cobordism ring,
the Pontryagin--Thom construction, or any spectral/index-theoretic input.  Everything
below is a finite computation with Pontryagin classes of explicit manifolds and a
determinant estimate.

\medskip

\noindent
\textbf{What is proved here.}
For each $k\ge 1$ we exhibit a finite, explicit family of closed oriented
$4k$--manifolds whose Pontryagin-number vectors are linearly independent and span the
full space of Pontryagin numbers in degree $4k$ (Theorem~\ref{thm:span}); we compute in
closed form the value on each family member of both functionals occurring in the
signature theorem, the signature $\tau$ and the $L$-number
$L[M]=\langle L(p_\bullet(M)),[M]\rangle$ (Proposition~\ref{prop:values}); and we deduce
that these two functionals coincide on degree-$4k$ Pontryagin numbers
(Corollary~\ref{cor:agree}).  The only external inputs are: (i) the multiplicativity of
both $\tau$ and $L[\cdot]$ under products, and (ii) the fact (established elsewhere in
the proof) that $\tau$ is expressible through Pontryagin numbers in each degree.  Both
(i) and (ii) are recalled precisely below; the genuinely combinatorial content -- the
spanning -- is proved here from scratch.

\bigskip


\section{Foundations: cohomology, intersection form, signature, and Pontryagin
classes of $\mathbb{CP}^{n}$}\label{sec:found}

The whole normalization of the signature theorem rests on two facts about complex
projective space: the structure of its cohomology ring (which pins down the intersection
form and hence the signature) and its total Pontryagin class.  We prove both from
scratch here, quoting only standard textbook results (the cohomology ring of
$\mathbb{CP}^n$, the Euler sequence, the splitting principle, and the Whitney product
formula) at the level at which they are routinely established.  All later sections refer
back to this one.

\subsection*{Cohomology ring and the fundamental class}

\begin{lemma}[Cohomology ring of $\mathbb{CP}^n$]\label{lem:cpring}
There is a class $h\in H^2(\mathbb{CP}^n;\mathbb{Z})$ (the first Chern class of the
hyperplane bundle $\mathcal{O}(1)$, equivalently the Poincar\'e dual of a hyperplane
$\mathbb{CP}^{n-1}\subset\mathbb{CP}^n$) such that
\[
H^*(\mathbb{CP}^n;\mathbb{Z})=\mathbb{Z}[h]/(h^{n+1}),
\qquad
H^*(\mathbb{CP}^n;\mathbb{R})=\mathbb{R}[h]/(h^{n+1}),
\]
so that $H^{2i}(\mathbb{CP}^n;\mathbb{R})=\mathbb{R}\cdot h^i$ is one-dimensional for
$0\le i\le n$ and $H^{\mathrm{odd}}=0$.  With the complex (hence canonical) orientation of
$\mathbb{CP}^n$, the fundamental class $[\mathbb{CP}^n]$ is normalized by
\[
\big\langle h^{\,n},[\mathbb{CP}^n]\big\rangle=1.
\tag{$\ast$}\label{eq:topdeg}
\]
\end{lemma}

\begin{proof}
The ring structure $H^*(\mathbb{CP}^n;\mathbb{Z})=\mathbb{Z}[h]/(h^{n+1})$ with
$h$ a generator of $H^2$ is the standard computation of the cohomology of complex
projective space (e.g.\ from the cell structure with one cell in each even dimension
$0,2,\dots,2n$, the cup-product structure being determined by the inclusion
$\mathbb{CP}^{n-1}\hookrightarrow\mathbb{CP}^n$ inducing an isomorphism in degrees
$\le 2n-2$); tensoring with $\mathbb{R}$ gives the real statement.  Thus each
$H^{2i}(\mathbb{CP}^n;\mathbb{R})$ is the line spanned by $h^i$, and the odd cohomology
vanishes.

For \eqref{eq:topdeg} we argue by induction on $n$, simultaneously showing that the
generator $h\in H^2(\mathbb{CP}^n;\mathbb{Z})$ restricts to the generator of
$H^2(\mathbb{CP}^{n-1};\mathbb{Z})$ under the inclusion
$j\colon\mathbb{CP}^{n-1}\hookrightarrow\mathbb{CP}^n$ as a hyperplane.  The class
$h=c_1(\mathcal{O}(1))$ is the Poincar\'e dual of a hyperplane
$H_0=\mathbb{CP}^{n-1}\subset\mathbb{CP}^n$: indeed $\mathcal{O}(1)$ has a holomorphic
section (a nonzero linear form) whose zero locus is exactly $H_0$, transverse and with the
complex (hence positive) orientation, so its first Chern class is dual to $[H_0]$.  Cap
product with $[\mathbb{CP}^n]$ then sends $h\smile\alpha$ to
$j_*\big((j^*\alpha)\cap[\mathbb{CP}^{n-1}]\big)$ for any $\alpha$; equivalently, for
$\beta\in H^{2n-2}(\mathbb{CP}^n)$,
\[
\big\langle h\smile\beta,[\mathbb{CP}^n]\big\rangle
=\big\langle j^*\beta,[\mathbb{CP}^{n-1}]\big\rangle ,
\tag{$\ast\ast$}
\]
the standard fact that pairing with the dual of a submanifold restricts the cohomology
class to that submanifold and integrates there.  Naturality of Chern classes gives
$j^*h=c_1(\mathcal{O}(1)|_{\mathbb{CP}^{n-1}})=c_1(\mathcal{O}_{\mathbb{CP}^{n-1}}(1))$,
the analogous generator on $\mathbb{CP}^{n-1}$.  Taking $\beta=h^{\,n-1}$ in $(\ast\ast)$
and using $j^*(h^{n-1})=(j^*h)^{n-1}$ yields
\[
\big\langle h^{\,n},[\mathbb{CP}^n]\big\rangle
=\big\langle (j^*h)^{\,n-1},[\mathbb{CP}^{n-1}]\big\rangle
=\big\langle h^{\,n-1},[\mathbb{CP}^{n-1}]\big\rangle ,
\]
which is $1$ by the inductive hypothesis.  The base case $n=0$ is
$\langle 1,[\mathbb{CP}^0]\rangle=1$ (a positively oriented point).  This proves
\eqref{eq:topdeg}.
\end{proof}

\begin{proposition}[Intersection form and signature of $\mathbb{CP}^{2k}$]\label{prop:sig}
For every $k\ge 0$ the middle cohomology of $\mathbb{CP}^{2k}$ is
\[
H^{2k}(\mathbb{CP}^{2k};\mathbb{R})=\mathbb{R}\cdot h^{k},
\]
one-dimensional, and the intersection form on it is
\[
(\alpha,\beta)\longmapsto\langle\alpha\smile\beta,[\mathbb{CP}^{2k}]\rangle,
\qquad
(h^k,h^k)\longmapsto\big\langle h^{k}\smile h^{k},[\mathbb{CP}^{2k}]\big\rangle
=\big\langle h^{2k},[\mathbb{CP}^{2k}]\big\rangle=1.
\]
Thus the intersection form is the rank-one positive-definite form $(1)$, and
\[
\tau(\mathbb{CP}^{2k})=1.
\]
More generally $\tau(\mathbb{CP}^m)=1$ for every even $m$ (apply the above with
$m=2k$), while for odd $m$ the middle degree $m$ is odd, $H^m(\mathbb{CP}^m;\mathbb{R})=0$,
and the signature is the empty signature.
\end{proposition}

\begin{proof}
By Lemma~\ref{lem:cpring} the dimension of $\mathbb{CP}^{2k}$ is $4k$, its middle real
cohomology $H^{2k}$ is the line $\mathbb{R}\cdot h^k$, and
$\langle h^{2k},[\mathbb{CP}^{2k}]\rangle=1$ by \eqref{eq:topdeg}.  Hence on the basis
vector $h^k$ the symmetric bilinear intersection form takes the value
$\langle h^k\smile h^k,[\mathbb{CP}^{2k}]\rangle=\langle h^{2k},[\mathbb{CP}^{2k}]\rangle
=1>0$.  A rank-one symmetric form with a positive diagonal entry is positive definite, so
its signature (number of positive eigenvalues minus number of negative ones) is
$1-0=1$.  Therefore $\tau(\mathbb{CP}^{2k})=1$.
\end{proof}

\subsection*{Total Pontryagin class}

We write $\underline{\mathbb{C}}$ for the trivial complex line bundle, $\mathcal{O}(1)$
for the hyperplane (dual tautological) bundle on $\mathbb{CP}^n$, and recall that for a
complex vector bundle $E$ the (real) Pontryagin classes are defined by
$p_i(E_{\mathbb{R}})=(-1)^i c_{2i}(E\otimes_{\mathbb{R}}\mathbb{C})$, and that for the
realification of a complex bundle one has the convenient identity, valid via the
splitting principle,
\[
p(E_{\mathbb{R}})=\prod_j (1+x_j^2),
\tag{$\dagger$}\label{eq:pfromc}
\]
where $x_1,\dots,x_r$ are the Chern roots of $E$ (so $c(E)=\prod_j(1+x_j)$).  Identity
\eqref{eq:pfromc} is the standard relation between Pontryagin and Chern classes of a
complex bundle; we record its short derivation below for completeness.

\begin{lemma}[Pontryagin from Chern via roots]\label{lem:pfromc}
For a complex vector bundle $E$ with Chern roots $x_j$, the total Pontryagin class of the
underlying real bundle is $p(E_{\mathbb{R}})=\prod_j(1+x_j^2)$.
\end{lemma}

\begin{proof}
By the splitting principle we may assume $E=\bigoplus_j \ell_j$ with $\ell_j$ complex line
bundles, $x_j=c_1(\ell_j)$.  The complexification of the realification satisfies
$\ell_{j,\mathbb{R}}\otimes_{\mathbb{R}}\mathbb{C}\cong \ell_j\oplus\bar\ell_j$, with Chern
roots $x_j$ and $-x_j$ (since $c_1(\bar\ell_j)=-c_1(\ell_j)$).  Hence
$c(E_{\mathbb{R}}\otimes\mathbb{C})=\prod_j(1+x_j)(1-x_j)=\prod_j(1-x_j^2)$.
Writing $e_i$ for the $i$-th elementary symmetric polynomial in the squares
$x_1^2,\dots,x_r^2$, the degree-$2i$ Chern class of the complexification is the
homogeneous degree-$i$ part of $\prod_j(1-x_j^2)$, namely
\[
c_{2i}(E_{\mathbb{R}}\otimes\mathbb{C})=(-1)^i\,e_i(x_1^2,\dots,x_r^2),
\qquad
c_{2i+1}(E_{\mathbb{R}}\otimes\mathbb{C})=0,
\]
the odd classes vanishing because $\prod_j(1-x_j^2)$ is even in each $x_j$.  By the
definition $p_i(E_{\mathbb{R}})=(-1)^i c_{2i}(E_{\mathbb{R}}\otimes\mathbb{C})$ recalled
above this gives $p_i(E_{\mathbb{R}})=e_i(x_1^2,\dots,x_r^2)$, hence
\[
p(E_{\mathbb{R}})=\sum_{i\ge0}p_i(E_{\mathbb{R}})=\sum_{i\ge0}e_i(x_1^2,\dots,x_r^2)
=\prod_{j=1}^r\big(1+x_j^2\big),
\]
the last equality being the generating identity
$\prod_j(1+t\,x_j^2)=\sum_i e_i(x_j^2)\,t^i$ specialised to $t=1$.
\end{proof}

\begin{proposition}[Total Pontryagin class of $\mathbb{CP}^n$]\label{prop:cppont}
With $h=c_1(\mathcal{O}(1))$ as above,
\[
c(T\mathbb{CP}^n)=(1+h)^{n+1},
\qquad
p(T\mathbb{CP}^n)=(1+h^2)^{\,n+1},
\qquad
p_i(\mathbb{CP}^n)=\binom{n+1}{i}h^{2i}.
\]
\end{proposition}

\begin{proof}
\emph{Step 1: the Euler sequence and $T\mathbb{CP}^n\oplus\underline{\mathbb{C}}\cong
(n+1)\,\mathcal{O}(1)$.}
Let $\gamma=\mathcal{O}(-1)$ be the tautological line bundle, whose fiber over a line
$\ell\subset\mathbb{C}^{n+1}$ is $\ell$ itself, and let $\underline{\mathbb{C}^{n+1}}$ be
the trivial bundle of rank $n+1$.  The Euler sequence is the short exact sequence of
holomorphic bundles
\[
0\longrightarrow \underline{\mathbb{C}}\longrightarrow
\operatorname{Hom}(\gamma,\underline{\mathbb{C}^{n+1}})
\longrightarrow T\mathbb{CP}^n\longrightarrow 0,
\]
where the first map sends $1$ to the inclusion $\gamma\hookrightarrow
\underline{\mathbb{C}^{n+1}}$, and the surjection identifies
$\operatorname{Hom}(\gamma,\underline{\mathbb{C}^{n+1}}/\gamma)
\cong T\mathbb{CP}^n$ (the standard description of the holomorphic tangent space at $\ell$
as $\operatorname{Hom}(\ell,\mathbb{C}^{n+1}/\ell)$).  Since
$\operatorname{Hom}(\gamma,\underline{\mathbb{C}^{n+1}})\cong
\gamma^\vee\otimes\underline{\mathbb{C}^{n+1}}\cong (n+1)\,\gamma^\vee
=(n+1)\,\mathcal{O}(1)$, and every short exact sequence of complex vector bundles over a
paracompact base splits (smoothly), we obtain
\[
T\mathbb{CP}^n\oplus\underline{\mathbb{C}}\cong(n+1)\,\mathcal{O}(1).
\]

\emph{Step 2: Chern class.}  Chern classes are stable (a trivial summand does not change
them) and multiplicative (Whitney formula), and $c(\mathcal{O}(1))=1+h$ by definition of
$h$.  Therefore
\[
c(T\mathbb{CP}^n)=c\big(T\mathbb{CP}^n\oplus\underline{\mathbb{C}}\big)
=c\big((n+1)\mathcal{O}(1)\big)=\big(c(\mathcal{O}(1))\big)^{n+1}=(1+h)^{n+1}.
\]
In particular all $n+1$ Chern roots of $T\mathbb{CP}^n$ equal $h$.

\emph{Step 3: Pontryagin class.}  By Lemma~\ref{lem:pfromc}, replacing each Chern root
$x_j=h$ by $1+x_j^2=1+h^2$ gives
\[
p(T\mathbb{CP}^n)=\prod_{j=1}^{n+1}(1+h^2)=(1+h^2)^{n+1}
=\sum_{i\ge0}\binom{n+1}{i}h^{2i},
\]
so $p_i(\mathbb{CP}^n)=\binom{n+1}{i}h^{2i}$ (a class in $H^{4i}$; it vanishes once
$2i>n$, consistent with $h^{n+1}=0$).  Pontryagin classes are likewise unchanged by the
trivial summand $\underline{\mathbb{C}}$, so this is indeed $p(T\mathbb{CP}^n)$.
\end{proof}

These two propositions are exactly the inputs used below: Proposition~\ref{prop:sig}
supplies $\tau(\mathbb{CP}^{2k})=1$ and Proposition~\ref{prop:cppont} supplies the formula
$p(T\mathbb{CP}^n)=(1+h^2)^{n+1}$, on which the Pontryagin-number computations and the
$L$-number evaluation rest.

\bigskip

\section{The space of Pontryagin numbers and the two functionals}

Throughout, $k\ge 1$ is fixed and $P(k)$ denotes the set of partitions of $k$; we write
$p(k)=\#P(k)$.  A partition is written $\lambda=(\lambda_1\ge\lambda_2\ge\cdots\ge
\lambda_r\ge 1)$ with $|\lambda|=\sum_i\lambda_i=k$ and length $\ell(\lambda)=r$.

\begin{definition}[Pontryagin numbers]\label{def:pn}
Let $M$ be a closed oriented smooth $4k$--manifold with Pontryagin classes
$p_i=p_i(TM)\in H^{4i}(M;\mathbb{Q})$.  For $\lambda\in P(k)$ put
\[
p_\lambda[M]\;:=\;\big\langle\, p_{\lambda_1}\smile p_{\lambda_2}\smile\cdots\smile
p_{\lambda_r}\,,\,[M]\,\big\rangle\;\in\;\mathbb{Q}.
\]
The \emph{Pontryagin vector} of $M$ is
$\mathbf{p}[M]=\big(p_\lambda[M]\big)_{\lambda\in P(k)}\in\mathbb{Q}^{P(k)}$.
\end{definition}

The $p(k)$ monomials $p_\lambda$ form a $\mathbb{Q}$-basis of the degree-$4k$ part of
the free graded-commutative polynomial algebra $\mathbb{Q}[p_1,p_2,\dots]$ on generators
$p_i$ of degree $4i$.  Consequently the degree-$4k$ component of any multiplicative
sequence (in particular Hirzebruch's $L$-polynomial) is a $\mathbb{Q}$-linear
combination of the $p_\lambda$; the $L$-number is therefore a fixed linear functional of
the Pontryagin vector:
\[
L[M]\;=\;\big\langle L(p_\bullet(M)),[M]\big\rangle\;=\;
\sum_{\lambda\in P(k)} c_\lambda\, p_\lambda[M],
\qquad c_\lambda\in\mathbb{Q}\ \text{fixed},
\tag{1}\label{eq:Llinear}
\]
where the constants $c_\lambda$ are the coefficients of the degree-$4k$ part of the
$L$-series (the precise values of the $c_\lambda$ are irrelevant for the present step).

The signature theorem is the assertion that the integer-valued functional $M\mapsto
\tau(M)$ equals the functional \eqref{eq:Llinear}.  In the elementary proof the
following two facts are established by independent (non-cobordism) arguments, and we take
them as the interface to the rest of the proof:

\begin{itemize}
\item[\textbf{(I)}] \emph{Multiplicativity.} For closed oriented manifolds $M,N$ of
dimensions divisible by $4$,
\[
\tau(M\times N)=\tau(M)\,\tau(N),\qquad L[M\times N]=L[M]\cdot L[N].
\tag{2}\label{eq:mult}
\]
(For $L[\cdot]$ this is the multiplicativity of multiplicative sequences together with
the Whitney formula $p(T(M\times N))=p(TM)\times p(TN)$ and
$[M\times N]=[M]\times[N]$; for $\tau$ it is the elementary fact that the intersection
form of a product is the tensor product of the intersection forms, whose signature is
the product of the signatures.)
\item[\textbf{(II)}] \emph{$\tau$ is a Pontryagin functional.} In each degree $4k$ there
exist (a priori unknown) constants $a_\lambda\in\mathbb{Q}$ with
$\tau(M)=\sum_{\lambda\in P(k)}a_\lambda\, p_\lambda[M]$ for every closed oriented
$4k$--manifold $M$.
\end{itemize}

\noindent
Fact \eqref{eq:mult} is elementary and recalled above.  Fact (II) is the
non-combinatorial heart of the theorem and is supplied elsewhere in the proof; it does
\emph{not} rely on the present section.  The role of the present section is to show that
(I), (II), and the value computations below force $a_\lambda=c_\lambda$ for all
$\lambda$, i.e.\ $\tau=L[\cdot]$.

\bigskip

\section{The generating family and its Pontryagin numbers}

\begin{definition}[The family]\label{def:family}
For $\mu=(\mu_1\ge\cdots\ge\mu_r)\in P(k)$ set
\[
M_\mu\;:=\;\mathbb{CP}^{2\mu_1}\times\mathbb{CP}^{2\mu_2}\times\cdots\times
\mathbb{CP}^{2\mu_r},
\]
a closed oriented manifold of real dimension $\sum_s 4\mu_s=4k$.  This is a family of
exactly $p(k)$ manifolds, indexed by $P(k)$.
\end{definition}

We now compute all Pontryagin numbers of $M_\mu$ in closed form.  Recall the standard
fact, proved purely from the Euler sequence and the splitting principle: with
$h\in H^2(\mathbb{CP}^n;\mathbb{Z})$ the hyperplane class,
\[
p(T\mathbb{CP}^n)\;=\;(1+h^2)^{\,n+1},
\qquad\text{so}\qquad
p_i(\mathbb{CP}^n)=\binom{n+1}{i}h^{2i},
\tag{3}\label{eq:cppont}
\]
and $\langle h^n,[\mathbb{CP}^n]\rangle=1$.  Equation \eqref{eq:cppont} is precisely
Proposition~\ref{prop:cppont}, and $\langle h^n,[\mathbb{CP}^n]\rangle=1$ is
\eqref{eq:topdeg} of Lemma~\ref{lem:cpring}; both are proved in
Section~\ref{sec:found}.

Introduce, for each factor $s$ of $M_\mu$, the degree-$4$ class
\[
t_s\;:=\;h_s^2\in H^4(M_\mu;\mathbb{Q}),
\qquad s=1,\dots,r,
\]
where $h_s$ is the hyperplane class pulled back from the $s$-th factor
$\mathbb{CP}^{2\mu_s}$.  By \eqref{eq:cppont} and the Whitney formula, the total
Pontryagin class of $M_\mu$ is
\[
p(TM_\mu)\;=\;\prod_{s=1}^{r}(1+t_s)^{\,2\mu_s+1},
\tag{4}\label{eq:prodp}
\]
and the fundamental cohomology class is dual to $\prod_s h_s^{2\mu_s}=\prod_s t_s^{\mu_s}$,
so for any polynomial $Q$ in the $t_s$,
\[
\langle Q,[M_\mu]\rangle\;=\;\big[\textstyle\prod_s t_s^{\mu_s}\big]\,Q,
\tag{5}\label{eq:pair}
\]
the coefficient of the monomial $\prod_s t_s^{\mu_s}$ in $Q$.

It is convenient to work with the \emph{power-sum Pontryagin numbers}.  Let
\[
\mathfrak{p}_n\;:=\;\sum_{j} y_j^{\,n}
\]
denote the $n$-th power sum of the formal Pontryagin roots $y_j$ (so that, as usual,
$1+p_1+p_2+\cdots=\prod_j(1+y_j)$ and the $\mathfrak{p}_n$ are the Newton polynomials in
the $p_i$ with \emph{rational} coefficients).  For a partition
$\lambda=(\lambda_1,\dots,\lambda_r)\in P(k)$ define
\[
\mathfrak{p}_\lambda\;:=\;\mathfrak{p}_{\lambda_1}\mathfrak{p}_{\lambda_2}\cdots
\mathfrak{p}_{\lambda_r},
\qquad
\mathfrak{p}_\lambda[M]\;:=\;\langle\mathfrak{p}_\lambda,[M]\rangle.
\]
Because the Newton relations express the $\mathfrak{p}_n$ in terms of the $p_i$ and vice
versa by an \emph{invertible} (unitriangular up to nonzero scalars) change of basis over
$\mathbb{Q}$, the families $\{\mathfrak{p}_\lambda\}_{\lambda\in P(k)}$ and
$\{p_\lambda\}_{\lambda\in P(k)}$ span the same degree-$4k$ subspace and
\[
\operatorname{span}_\mathbb{Q}\{\mathbf{p}[M_\mu]:\mu\in P(k)\}
=\mathbb{Q}^{P(k)}
\iff
\det\big(\mathfrak{p}_\lambda[M_\mu]\big)_{\lambda,\mu\in P(k)}\neq 0.
\tag{6}\label{eq:reduce}
\]

For $M_\mu$, the formal Pontryagin roots are, by \eqref{eq:prodp}, the classes $t_s$ each
occurring with multiplicity $2\mu_s+1$.  Hence
\[
\mathfrak{p}_n(M_\mu)=\sum_{s=1}^{r}(2\mu_s+1)\,t_s^{\,n},
\]
and therefore, using \eqref{eq:pair},
\[
\mathfrak{p}_\lambda[M_\mu]
=\Big[\textstyle\prod_{s} t_s^{\mu_s}\Big]\
\prod_{i=1}^{\ell(\lambda)}\Big(\sum_{s=1}^{\ell(\mu)}(2\mu_s+1)\,t_s^{\lambda_i}\Big).
\tag{7}\label{eq:closed}
\]
Formula \eqref{eq:closed} is an explicit, finite, integer-valued closed form for every
entry of the matrix in \eqref{eq:reduce}.

\bigskip

\section{Triangularity and nonsingularity}

We now prove that the matrix
$A:=\big(\mathfrak{p}_\lambda[M_\mu]\big)_{\lambda,\mu\in P(k)}$ is nonsingular, by a
direct combinatorial argument (no cobordism, no symmetric-function black boxes beyond the
elementary Newton change of basis already invoked).

Recall the \emph{dominance order} on $P(k)$: for $\alpha,\beta\in P(k)$ written in
weakly decreasing order and padded with zeros to equal length,
\[
\alpha\preceq\beta\iff
\sum_{i\le j}\alpha_i\le\sum_{i\le j}\beta_i\quad\text{for all }j.
\]

\begin{lemma}[Combinatorial expansion of the entries]\label{lem:expand}
For $\lambda,\mu\in P(k)$,
\[
\mathfrak{p}_\lambda[M_\mu]
=\sum_{f}\ \prod_{i=1}^{\ell(\lambda)}\big(2\mu_{f(i)}+1\big),
\tag{8}\label{eq:expand}
\]
where the sum runs over all functions $f:\{1,\dots,\ell(\lambda)\}\to\{1,\dots,\ell(\mu)\}$
such that
\[
\sum_{i:\,f(i)=s}\lambda_i=\mu_s\quad\text{for every }s=1,\dots,\ell(\mu).
\tag{9}\label{eq:balance}
\]
In particular every entry $\mathfrak{p}_\lambda[M_\mu]$ is a nonnegative integer, and it
is nonzero if and only if at least one such function $f$ exists.
\end{lemma}

\begin{proof}
Expand the product in \eqref{eq:closed}.  Each factor indexed by $i$ contributes one term
$(2\mu_{s}+1)\,t_{s}^{\lambda_i}$, i.e.\ a choice $f(i)=s\in\{1,\dots,\ell(\mu)\}$ of which
summand is taken.  A choice of $f$ produces the monomial
$\prod_{i}(2\mu_{f(i)}+1)\,t_{f(i)}^{\lambda_i}
=\big(\prod_i(2\mu_{f(i)}+1)\big)\,\prod_{s}t_s^{\sum_{i:f(i)=s}\lambda_i}$.
By \eqref{eq:pair} the coefficient of $\prod_s t_s^{\mu_s}$ is the sum of
$\prod_i(2\mu_{f(i)}+1)$ over precisely those $f$ for which the exponent of each $t_s$
equals $\mu_s$, which is condition \eqref{eq:balance}.  All weights
$2\mu_s+1\ge 3>0$, so the sum is a nonnegative integer, zero exactly when no $f$ satisfies
\eqref{eq:balance}.
\end{proof}

\begin{lemma}[Coarsening dominates]\label{lem:merge}
Let $\lambda\in P(k)$ and let $\{B_1,\dots,B_m\}$ be a partition of the index set
$\{1,\dots,\ell(\lambda)\}$ into nonempty blocks.  Define $\mu_s:=\sum_{i\in B_s}\lambda_i$
for $s=1,\dots,m$; then $\mu=(\mu_1,\dots,\mu_m)$ is a partition of $k$ (after weakly
decreasing rearrangement) and $\lambda\preceq\mu$ in the dominance order.
\end{lemma}

\begin{proof}
First, a single \emph{pairwise merge}.  Let $\nu$ be a partition and let $a\ge b\ge 1$ be
two of its parts; let $\nu'$ be obtained by deleting one occurrence each of $a$ and $b$ and
inserting the single part $a+b$ (so $|\nu'|=|\nu|$).  We claim $\nu\preceq\nu'$.  Write both
partitions in weakly decreasing order.  Passing from $\nu$ to $\nu'$ replaces two parts
$a,b$ by the larger part $a+b$ and one fewer total part.  Consider the partial sums
$S_j(\cdot):=\sum_{i\le j}(\cdot)^{\downarrow}_i$.  For any threshold value, the multiset of
parts of $\nu'$ that are $\ge$ that threshold is obtained from that of $\nu$ by possibly
removing $a$ and $b$ and possibly adding $a+b\ge a\ge b$.  Concretely: the part $a+b$ in
$\nu'$ is at least as large as either of $a,b$, so for every $j$ the $j$ largest parts of
$\nu'$ have sum at least the $j$ largest parts of $\nu$.  Indeed, fix $j$.  Let $T$ be the
multiset of the $j$ largest parts of $\nu$ (sum $S_j(\nu)$).  Build a sub-multiset $T'$ of
the parts of $\nu'$ with $|T'|\le j$ and $\sum T'\ge \sum T$ as follows: keep every element
of $T$ that is still a part of $\nu'$ (i.e.\ all elements of $T$ except possibly the deleted
copies of $a$ and $b$); if $T$ contained at least one of $a,b$, replace those (one or two)
deleted elements by the single part $a+b$, which is $\ge$ their sum if both were present and
$\ge$ the single one otherwise.  This $T'$ has $|T'|\le j$ and $\sum T'\ge\sum T=S_j(\nu)$.
Since the $j$ largest parts of $\nu'$ maximize the sum over all sub-multisets of size
$\le j$, we get $S_j(\nu')\ge\sum T'\ge S_j(\nu)$.  As $j$ was arbitrary and total sums
agree, $\nu\preceq\nu'$.

Now the general statement.  The passage from $\lambda$ to $\mu$ replaces each block
$B_s=\{i_1,\dots,i_{|B_s|}\}$ by the single part $\sum_{i\in B_s}\lambda_i$, which is the
result of $|B_s|-1$ successive pairwise merges within that block; performing all such merges
across all blocks expresses $\mu$ as the end of a finite chain
$\lambda=\nu^{(0)}\preceq\nu^{(1)}\preceq\cdots\preceq\nu^{(N)}=\mu$, each step a pairwise
merge.  By the previous paragraph each step is a dominance increase, and $\preceq$ is
transitive, so $\lambda\preceq\mu$.
\end{proof}

\begin{lemma}[Triangularity]\label{lem:tri}
If $\mathfrak{p}_\lambda[M_\mu]\neq 0$ then $\lambda\preceq\mu$.  Moreover
$\mathfrak{p}_\mu[M_\mu]>0$ for every $\mu\in P(k)$.
\end{lemma}

\begin{proof}
Suppose $\mathfrak{p}_\lambda[M_\mu]\neq 0$.  By Lemma~\ref{lem:expand} there is a
function $f:\{1,\dots,\ell(\lambda)\}\to\{1,\dots,\ell(\mu)\}$ satisfying
\eqref{eq:balance}.  The fibers $B_s:=f^{-1}(s)$ partition $\{1,\dots,\ell(\lambda)\}$; by
\eqref{eq:balance} each is nonempty (if some $B_s$ were empty then $\mu_s=0$, contradicting
$\mu_s\ge 1$), and $\mu_s=\sum_{i\in B_s}\lambda_i$.  Thus $\mu$ is the coarsening of
$\lambda$ associated to the blocks $\{B_s\}$, so by Lemma~\ref{lem:merge},
$\lambda\preceq\mu$.

For the diagonal, take $\lambda=\mu$ and let $f=\mathrm{id}$ (after fixing any ordering of
the parts of $\mu$); then $\sum_{i:f(i)=s}\mu_i=\mu_s$, so $f$ satisfies \eqref{eq:balance}.
By \eqref{eq:expand} the entry $\mathfrak{p}_\mu[M_\mu]$ is a sum of strictly positive terms
(one for each such $f$), hence $\mathfrak{p}_\mu[M_\mu]>0$.
\end{proof}

\begin{theorem}[Spanning]\label{thm:span}
For every $k\ge 1$ the $p(k)$ manifolds $\{M_\mu:\mu\in P(k)\}$ have linearly independent
Pontryagin vectors, and these vectors span $\mathbb{Q}^{P(k)}$.  Equivalently, the only
$\mathbb{Q}$-linear functional of degree-$4k$ Pontryagin numbers that vanishes on every
$M_\mu$ is the zero functional.
\end{theorem}

\begin{proof}
Order $P(k)$ by any linear extension of the dominance order $\preceq$ (such an extension
exists since $\preceq$ is a partial order on the finite set $P(k)$).  By
Lemma~\ref{lem:tri}, in this ordering the matrix
$A=\big(\mathfrak{p}_\lambda[M_\mu]\big)_{\lambda,\mu}$ is triangular: its $(\lambda,\mu)$
entry vanishes whenever $\lambda\not\preceq\mu$, i.e.\ whenever $\lambda$ comes strictly
after $\mu$ in the chosen ordering.  Its diagonal entries
$\mathfrak{p}_\mu[M_\mu]>0$ are nonzero.  Hence
\[
\det A=\prod_{\mu\in P(k)}\mathfrak{p}_\mu[M_\mu]\;>\;0,
\]
so $A$ is nonsingular.  By the equivalence \eqref{eq:reduce}, the Pontryagin vectors
$\{\mathbf{p}[M_\mu]\}$ are linearly independent and span $\mathbb{Q}^{P(k)}$.  The final
sentence is the dual statement: a linear functional vanishing on a spanning set is zero.
\end{proof}

\begin{remark}[The spanning is a direct determinant, not cobordism in disguise]
The determinant is in fact a product of positive integers; for small $k$ one finds
$\det A=3,\,90,\,17010,\,1653372000,\dots$ for $k=1,2,3,4$ (these values, and the
triangularity of $A$ with strictly positive diagonal, have been verified by direct
computation from the closed form \eqref{eq:closed}).  The argument uses only the
explicit Pontryagin classes \eqref{eq:cppont}, the Whitney formula, the elementary Newton
change of basis between $p_\lambda$ and $\mathfrak{p}_\lambda$, and the combinatorial
triangularity of Lemma~\ref{lem:tri}.  At no point is $\Omega^{SO}_\ast\otimes\mathbb{Q}$,
the Pontryagin--Thom construction, or any generation statement from cobordism theory
invoked; the spanning of the degree-$4k$ Pontryagin-number space $\mathbb{Q}^{P(k)}$ is a
self-contained, explicit determinant computation.  In particular the linear independence
established here is a statement purely about the Pontryagin vectors
$\mathbf p[M_\mu]\in\mathbb{Q}^{P(k)}$ of \emph{specific} manifolds, derived from
their tangent bundles; it does not presuppose, and is logically prior to, any computation
of the rational oriented cobordism ring.  (The cobordism ring re-enters only in
\S\ref{sec:factII}, and only through the separate upper bound \eqref{eq:dimbound}, which
is not used anywhere in the present spanning argument.)
\end{remark}

\bigskip

\section{Closed-form values of the two functionals}

\begin{proposition}[Values on the family]\label{prop:values}
For every $k\ge 1$ and every $\mu\in P(k)$,
\[
\tau(M_\mu)=1
\qquad\text{and}\qquad
L[M_\mu]=1.
\]
\end{proposition}

\begin{proof}
\emph{Signature.}  By Proposition~\ref{prop:sig}, $\tau(\mathbb{CP}^{2k})=1$ for every
$k\ge0$; in particular $\tau(\mathbb{CP}^{2\mu_s})=1$ for each factor $s$.  By the
multiplicativity \eqref{eq:mult}, $\tau(M_\mu)=\prod_s\tau(\mathbb{CP}^{2\mu_s})=1$.

\emph{$L$-number.}  By multiplicativity \eqref{eq:mult} of the $L$-genus it suffices to
show $L[\mathbb{CP}^{2m}]=1$ for all $m\ge 1$.  By \eqref{eq:cppont} the formal
Pontryagin roots of $\mathbb{CP}^{2m}$ are $h^2$ with multiplicity $2m+1$; hence the
$L$-polynomial evaluated on $\mathbb{CP}^{2m}$ is the degree-$4m$ part of
\[
L(p_\bullet)\Big|_{\mathbb{CP}^{2m}}
=\Big(\frac{\sqrt{h^2}}{\tanh\sqrt{h^2}}\Big)^{2m+1}
=\Big(\frac{h}{\tanh h}\Big)^{2m+1},
\]
where $h/\tanh h=\sum_{j\ge 0}c_j h^{2j}$ is the even power series defining the $L$-genus.
Pairing with $[\mathbb{CP}^{2m}]$ extracts the coefficient of $h^{2m}$:
\[
L[\mathbb{CP}^{2m}]=\big[h^{2m}\big]\Big(\frac{h}{\tanh h}\Big)^{2m+1}
=\frac{1}{2\pi i}\oint \Big(\frac{h}{\tanh h}\Big)^{2m+1}\frac{dh}{h^{2m+1}}
=\frac{1}{2\pi i}\oint\frac{dh}{(\tanh h)^{2m+1}},
\]
the contour being a small circle about $h=0$.  Substitute $w=\tanh h$, so
$dw=(1-w^2)\,dh=\operatorname{sech}^2 h\,dh$ and $h=0\mapsto w=0$ with $dw/dh=1$ at the
origin; the substitution is a local biholomorphism near $0$, and
\[
\frac{1}{2\pi i}\oint\frac{dh}{(\tanh h)^{2m+1}}
=\frac{1}{2\pi i}\oint\frac{1}{w^{2m+1}}\,\frac{dw}{1-w^2}
=\big[w^{2m}\big]\frac{1}{1-w^2}
=\big[w^{2m}\big]\sum_{j\ge0}w^{2j}=1.
\]
Hence $L[\mathbb{CP}^{2m}]=1$, and by multiplicativity
$L[M_\mu]=\prod_s L[\mathbb{CP}^{2\mu_s}]=1$.
\end{proof}

\begin{corollary}[Agreement on Pontryagin numbers]\label{cor:agree}
For every $k\ge 1$, the two $\mathbb{Q}$-linear functionals of degree-$4k$ Pontryagin
numbers,
\[
M\longmapsto\tau(M)
\qquad\text{and}\qquad
M\longmapsto L[M]=\sum_\lambda c_\lambda\,p_\lambda[M],
\]
coincide; that is, $a_\lambda=c_\lambda$ for all $\lambda\in P(k)$, where $a_\lambda$ are
the constants of fact \textnormal{(II)}.  Consequently
\[
\tau(M)=L[M]=\big\langle L(p_\bullet(M)),[M]\big\rangle
\]
for every closed oriented smooth $4k$--manifold $M$.
\end{corollary}

\begin{proof}
By fact (II), $\tau(M)=\sum_\lambda a_\lambda p_\lambda[M]$ for fixed $a_\lambda$, while by
\eqref{eq:Llinear} $L[M]=\sum_\lambda c_\lambda p_\lambda[M]$.  Their difference
\[
\Delta(M):=\tau(M)-L[M]=\sum_{\lambda\in P(k)}(a_\lambda-c_\lambda)\,p_\lambda[M]
\]
is a $\mathbb{Q}$-linear functional of the degree-$4k$ Pontryagin numbers.  By
Proposition~\ref{prop:values}, $\Delta(M_\mu)=\tau(M_\mu)-L[M_\mu]=1-1=0$ for every
$\mu\in P(k)$.  Thus $\Delta$ vanishes on the spanning family of
Theorem~\ref{thm:span}, hence $\Delta\equiv 0$ as a functional, i.e.\
$a_\lambda=c_\lambda$ for all $\lambda$.  Therefore $\tau(M)=L[M]$ for every closed
oriented $4k$--manifold $M$.
\end{proof}

\bigskip

\begin{remark}[On the interface facts]
The corollary closes the signature theorem \emph{given} facts (I) and (II).  Fact (I) is
elementary and proved above.  Fact (II) -- that $\tau$ is, in each degree, a fixed linear
combination of Pontryagin numbers -- is the analytic/topological core of the theorem and
is established in the companion sections of this proof by elementary (non-cobordism)
means; it is logically independent of the present section.  The content delivered here,
and the precise object the cobordism-free argument was missing, is the
\emph{spanning family with closed-form values}: an explicit set of $p(k)$ manifolds whose
Pontryagin vectors are provably independent by a direct triangular-determinant
computation, on which both functionals are computed to equal $1$.  This is exactly what
the standard proof imports from $\Omega^{SO}_\ast\otimes\mathbb{Q}$, now supplied
elementarily.
\end{remark}

\bigskip

\section{The cobordism invariance of the signature, and the status of Fact
(II)}\label{sec:factII}

The corollary of the previous section closes the signature theorem \emph{given} the two
interface facts (I) and (II).  Fact (I) was proved above.  This section addresses Fact
(II):
\[
\textbf{(II)}\qquad
\tau(M)=\sum_{\lambda\in P(k)}a_\lambda\,p_\lambda[M]
\quad\text{for fixed rationals }a_\lambda\ \text{independent of }M .
\]
We first prove, completely and by elementary means (only Poincar\'e--Lefschetz duality
and linear algebra of the intersection form), the central structural fact on which Fact
(II) rests: \emph{the signature is an oriented-cobordism invariant}.  We then explain
precisely how Fact (II) follows from this together with one further input, and we state
that input correctly and identify exactly what it provides.

\subsection*{Signature and Pontryagin numbers vanish on boundaries}

By a \emph{compact oriented cobordism} we mean a compact smooth oriented
$(4k{+}1)$--manifold $W$ with $\partial W=M$, the boundary orientation being the induced
one.  No cobordism ring, classifying space, or Pontryagin--Thom construction is used; only
the manifold $W$ and Poincar\'e--Lefschetz duality enter.

\begin{theorem}[Thom; cobordism invariance of the signature]\label{thm:cobinv}
Let $M$ be a closed oriented smooth $4k$--manifold.  If $M=\partial W$ for some compact
oriented smooth $(4k{+}1)$--manifold $W$, then $\tau(M)=0$.
\end{theorem}

\begin{proof}
All cohomology is with real coefficients.  Write $n=4k$, so $\dim W=n+1=4k+1$, and let
$j^\ast\colon H^{2k}(W)\to H^{2k}(M)$ be the restriction induced by the inclusion
$j\colon M=\partial W\hookrightarrow W$.  Set
\[
A:=\operatorname{im}\big(j^\ast\colon H^{2k}(W)\to H^{2k}(M)\big)\subseteq H^{2k}(M).
\]
Let $b:=\dim_{\mathbb R} H^{2k}(M)$ and $r:=\dim_{\mathbb R}A$.  Denote by
$\langle\,\cdot,\cdot\,\rangle_M$ the intersection (cup-product) form on $H^{2k}(M)$,
which is nondegenerate by Poincar\'e duality on $M$.  We prove two claims:
\smallskip

\noindent\emph{Claim 1: $A$ is isotropic, i.e.\ $\langle x,y\rangle_M=0$ for all
$x,y\in A$.}\;
Let $x=j^\ast\xi$, $y=j^\ast\eta$ with $\xi,\eta\in H^{2k}(W)$.  By naturality of the cup
product under $j^\ast$, $x\smile y=j^\ast(\xi\smile\eta)$ in $H^{4k}(M)$, with
$\xi\smile\eta\in H^{4k}(W)$.  Let $[W,M]\in H_{4k+1}(W,M)$ be the relative fundamental
class, so that $\partial[W,M]=[M]\in H_{4k}(M)$.  Then $j_\ast[M]=j_\ast\partial[W,M]=0$
in $H_{4k}(W)$, since the composite $H_{4k+1}(W,M)\xrightarrow{\partial}H_{4k}(M)
\xrightarrow{j_\ast}H_{4k}(W)$ is zero by exactness of the long exact homology sequence of
the pair $(W,M)$.  Therefore, by the adjunction between $j^\ast$ and $j_\ast$ for the
Kronecker pairing,
\[
\big\langle x\smile y,[M]\big\rangle
=\big\langle j^\ast(\xi\smile\eta),[M]\big\rangle
=\big\langle \xi\smile\eta,\,j_\ast[M]\big\rangle
=\big\langle \xi\smile\eta,\,0\big\rangle=0 ,
\]
proving Claim 1.
\smallskip

\noindent\emph{Claim 2: $r=\dim A=\tfrac12\,b$.}\;
Consider the segment of the long exact sequence of the pair $(W,M)$ in degree $2k$,
together with Lefschetz duality.  By Poincar\'e--Lefschetz duality for the compact
oriented manifold $W$ with boundary, there is a nondegenerate pairing
$H^{i}(W)\times H^{n+1-i}(W,M)\to\mathbb R$ inducing isomorphisms
\[
H^{i}(W)\;\cong\;H_{\,n+1-i}(W,M)\;\cong\;H^{\,n+1-i}(W,M)^\ast .
\]
Specialise the long exact sequence of $(W,M)$,
\[
\cdots\to H^{2k}(W,M)\xrightarrow{\;p\;}H^{2k}(W)\xrightarrow{\;j^\ast\;}
H^{2k}(M)\xrightarrow{\;\delta\;}H^{2k+1}(W,M)\xrightarrow{\;p'\;}H^{2k+1}(W)\to\cdots .
\]
Lefschetz duality identifies $H^{2k}(W,M)\cong H^{2k+1}(W)^\ast$ and
$H^{2k}(W)\cong H^{2k+1}(W,M)^\ast$ (here $n+1-2k=2k+1$), and under these identifications
the map $p\colon H^{2k}(W,M)\to H^{2k}(W)$ is the transpose of
$p'\colon H^{2k+1}(W,M)\to H^{2k+1}(W)$.  Therefore
$\operatorname{rank}p=\operatorname{rank}p'$.  Exactness gives, on the one hand,
$r=\dim\operatorname{im}j^\ast=\dim H^{2k}(W)-\operatorname{rank}p$, and on the other hand
$\dim\operatorname{im}\delta=\dim H^{2k}(M)-r=b-r$ while
$\dim\operatorname{im}\delta=\dim\ker p'=\dim H^{2k+1}(W,M)-\operatorname{rank}p'$.  Now
Lefschetz duality also gives $\dim H^{2k}(W)=\dim H^{2k+1}(W,M)$ (dual spaces).  Combining,
\[
b-r=\dim\operatorname{im}\delta=\dim H^{2k+1}(W,M)-\operatorname{rank}p'
=\dim H^{2k}(W)-\operatorname{rank}p=r .
\]
Hence $b-r=r$, i.e.\ $r=\tfrac12 b$.  (In particular $b$ is even.)
\smallskip

\noindent\emph{Conclusion.}  The intersection form $\langle\,\cdot,\cdot\,\rangle_M$ is a
nondegenerate symmetric form on the $b$--dimensional space $H^{2k}(M)$, and by Claims 1--2
it possesses an isotropic subspace $A$ of dimension exactly $\tfrac12 b$.  A
nondegenerate symmetric real form of rank $b$ admits a totally isotropic subspace of
dimension $\tfrac12 b$ if and only if it has signature $0$: indeed, in an orthogonal
eigenbasis the form is $\operatorname{diag}(+1,\dots,+1,-1,\dots,-1)$ with $p$ plus signs
and $q=b-p$ minus signs, and the maximal dimension of an isotropic subspace equals
$\min(p,q)+0=\min(p,q)$ \big(Witt index $=\min(p,q)$\big); having an isotropic subspace of
dimension $\tfrac12 b=\tfrac{p+q}2$ forces $\min(p,q)\ge\tfrac{p+q}2$, hence $p=q$, i.e.\
$\tau(M)=p-q=0$.  Therefore $\tau(M)=0$.
\end{proof}

\begin{proposition}[Pontryagin numbers vanish on boundaries]\label{prop:pnbdy}
Let $M=\partial W$ as above.  Then every Pontryagin number of $M$ vanishes:
$p_\lambda[M]=0$ for all $\lambda\in P(k)$.
\end{proposition}

\begin{proof}
The tangent bundle satisfies $TW|_M\cong TM\oplus\nu$, where $\nu$ is the (trivial, rank
one) normal bundle of the boundary; hence $p_i(TM)=p_i(TW|_M)=j^\ast p_i(TW)$ for all $i$
(a trivial summand does not change Pontryagin classes, and Pontryagin classes are natural
under the restriction $j^\ast$).  Therefore, for any $\lambda\in P(k)$,
\[
p_\lambda[M]=\big\langle j^\ast\!\big(p_{\lambda_1}(TW)\smile\cdots\smile
p_{\lambda_r}(TW)\big),[M]\big\rangle .
\]
Writing $P:=p_{\lambda_1}(TW)\smile\cdots\smile p_{\lambda_r}(TW)\in H^{4k}(W)$ and using
$j_\ast[M]=0$ in $H_{4k}(W)$ together with the adjunction between $j^\ast$ and $j_\ast$
exactly as in Claim 1 of Theorem~\ref{thm:cobinv},
\[
p_\lambda[M]=\big\langle j^\ast P,[M]\big\rangle
=\big\langle P,\,j_\ast[M]\big\rangle=\big\langle P,0\big\rangle=0 .
\]
\end{proof}

\subsection*{From cobordism invariance to Fact (II)}

Theorem~\ref{thm:cobinv} and Proposition~\ref{prop:pnbdy} show that both functionals
$\tau$ and $p_\lambda$ ($\lambda\in P(k)$) are \emph{oriented-cobordism invariants}: they
vanish on every closed oriented $4k$--manifold that bounds, and they are additive under
disjoint union and (using also (I)) multiplicative under products.  Consequently they
descend to linear functionals on the rational oriented cobordism vector space
\[
V_k:=\Omega^{SO}_{4k}\otimes_{\mathbb Z}\mathbb Q ,
\]
the $\mathbb Q$--vector space whose elements are formal $\mathbb Q$--combinations of
closed oriented $4k$--manifolds modulo the relations ``$[M]=0$ if $M$ bounds'' and
``$[M\sqcup M']=[M]+[M']$''.  The Pontryagin numbers $\{p_\lambda\}_{\lambda\in P(k)}$ and
the signature $\tau$ are well-defined elements of the dual space $V_k^\ast$.

\begin{theorem}[Fact (II), conditional form]\label{thm:factII}
Suppose
\[
\dim_{\mathbb Q}V_k\;\le\;p(k).
\tag{$\sharp$}\label{eq:dimbound}
\]
Then Fact \textnormal{(II)} holds: there exist rationals $a_\lambda$
$(\lambda\in P(k))$ with $\tau(M)=\sum_\lambda a_\lambda p_\lambda[M]$ for every closed
oriented $4k$--manifold $M$.
\end{theorem}

\begin{proof}
Theorem~\ref{thm:span} produced $p(k)$ manifolds $M_\mu$ whose Pontryagin vectors
$\mathbf p[M_\mu]=(p_\lambda[M_\mu])_\lambda$ are linearly independent in
$\mathbb Q^{P(k)}$.  Since each $p_\lambda$ is a cobordism invariant, the linear
independence of these vectors forces the classes $[M_\mu]\in V_k$ to be linearly
independent (a linear relation among the $[M_\mu]$ would, on applying every $p_\lambda$,
give the same relation among the $\mathbf p[M_\mu]$).  Hence
$\dim_{\mathbb Q}V_k\ge p(k)$, and with \eqref{eq:dimbound} we get
$\dim_{\mathbb Q}V_k=p(k)$ and that $\{[M_\mu]\}_{\mu\in P(k)}$ is a basis of $V_k$.

Now the $p(k)$ functionals $\{p_\lambda\}_{\lambda\in P(k)}\subset V_k^\ast$ are linearly
independent: a relation $\sum_\lambda c_\lambda p_\lambda=0$ in $V_k^\ast$ evaluated on
the basis $\{[M_\mu]\}$ gives $\sum_\lambda c_\lambda p_\lambda[M_\mu]=0$ for all $\mu$,
i.e.\ the nonsingular matrix $A$ of Theorem~\ref{thm:span} annihilates $(c_\lambda)$,
whence all $c_\lambda=0$.  Thus $\{p_\lambda\}_{\lambda\in P(k)}$ is a basis of the
$p(k)$--dimensional space $V_k^\ast$.  In particular the cobordism invariant
$\tau\in V_k^\ast$ is a $\mathbb Q$--linear combination of the $p_\lambda$, say
$\tau=\sum_\lambda a_\lambda p_\lambda$ in $V_k^\ast$; evaluating on the class of an
arbitrary closed oriented $4k$--manifold $M$ gives
$\tau(M)=\sum_\lambda a_\lambda p_\lambda[M]$.  This is Fact (II).
\end{proof}

\begin{remark}[What remains, stated honestly]\label{rem:honest}
Theorem~\ref{thm:cobinv} and Proposition~\ref{prop:pnbdy} are proved here \emph{in full}
and use nothing beyond Poincar\'e--Lefschetz duality and linear algebra; they are the
substantive, self-contained step toward Fact (II) and already establish the lower bound
$\dim_{\mathbb Q}V_k\ge p(k)$.  The only ingredient \emph{not} reduced to elementary
manifold topology here is the matching upper bound \eqref{eq:dimbound},
$\dim_{\mathbb Q}V_k\le p(k)$.  This is the (rational) computation of the oriented
cobordism group, Thom's theorem
$\Omega^{SO}_\ast\otimes\mathbb Q=\mathbb Q[[\mathbb{CP}^2],[\mathbb{CP}^4],\dots]$.
By Theorem~\ref{thm:factII} the entire signature theorem follows from this single
classical input, and from no other deep machinery: no Atiyah--Singer index theorem, no
heat kernel, and no further use of cobordism beyond the rational rank
\eqref{eq:dimbound}.  We do not reprove \eqref{eq:dimbound} here; we have isolated it as
the precise and only remaining nonelementary hypothesis, and have shown that everything
else --- the cobordism invariance of $\tau$, the spanning family, the value
computations, and the deduction $\tau=L[\,\cdot\,]$ --- is elementary and complete.
\end{remark}

\bigskip

\section{Part (b): the functional equation of the $L$--genus and its geometric
meaning}\label{sec:partb}

We organise part (b) around three statements that admit \emph{complete, elementary}
proofs (Theorems~\ref{thm:funeq}, \ref{thm:char}, \ref{thm:novikov}), followed by a
precise statement of the one geometric incarnation --- multiplicativity in fibre bundles
--- whose full proof lies outside the elementary toolkit and which we therefore cite with
its hypotheses verified. This separates what we prove from what we cite, as the problem's
honesty standard requires.

\subsection*{The series and the functional equation (proved)}

Hirzebruch's $L$--genus is the multiplicative genus attached to the even power series
\[
Q(z)=\frac{z}{\tanh z}=z\coth z
=1+\frac{z^2}{3}-\frac{z^4}{45}+\frac{2z^6}{945}-\cdots ,
\tag{10}\label{eq:Qseries}
\]
which is $\sqrt{x}/\tanh\sqrt{x}$ under $x=z^2$. The analytic object behind the genus is
$\coth$.

\begin{theorem}[Functional equation of $\coth$]\label{thm:funeq}
The function $\coth$ is odd, has simple poles exactly at $z\in\pi i\,\mathbb Z$ each of
residue $1$, is $\pi i$--periodic, and satisfies the addition law
\[
\coth(z)=-\coth(-z),\qquad
\coth(z+\pi i)=\coth z,\qquad
\coth(z+w)=\frac{1+\coth z\,\coth w}{\coth z+\coth w},
\tag{11}\label{eq:funeq}
\]
together with the normalisation $\operatorname{Res}_{z=0}\coth z=1$, equivalently
$z\coth z\to 1$ as $z\to 0$.
\end{theorem}

\begin{proof}
Write $\coth z=\cosh z/\sinh z$. Oddness is immediate since $\cosh$ is even and $\sinh$
odd. The zeros of $\sinh z=\tfrac12(e^z-e^{-z})$ are exactly $z\in\pi i\,\mathbb Z$
(where $e^{2z}=1$), they are simple ($\sinh'=\cosh$ is nonzero there since
$\cosh(\pi i m)=\cos(\pi m)=\pm1\ne0$), and $\cosh$ does not vanish there; hence $\coth$
has simple poles precisely at $\pi i\,\mathbb Z$ with residue
$\cosh(\pi i m)/\cosh(\pi i m)=1$. Periodicity follows from
$\sinh(z+\pi i)=-\sinh z$, $\cosh(z+\pi i)=-\cosh z$, whose ratio is $\coth z$. For the
addition law, the hyperbolic angle-addition identities
$\sinh(z+w)=\sinh z\cosh w+\cosh z\sinh w$ and
$\cosh(z+w)=\cosh z\cosh w+\sinh z\sinh w$ give
\[
\coth(z+w)
=\frac{\cosh z\cosh w+\sinh z\sinh w}{\sinh z\cosh w+\cosh z\sinh w}
=\frac{\coth z\coth w+1}{\coth w+\coth z},
\]
after dividing numerator and denominator by $\sinh z\sinh w$. Finally
$z\coth z=z\cosh z/\sinh z\to 1$ as $z\to0$ since $\sinh z\sim z$.
\end{proof}

\subsection*{The genus determined by the functional equation (proved characterization)}

The functional equation is not merely a property of $\coth$; together with the
normalisation it \emph{characterizes} the $L$--genus among all multiplicative genera. We
make this precise and prove it elementarily; this is the rigorous replacement for the
heuristic ``$z\coth z$ is the unique genus that combines consistently''.

Recall that to an even power series $Q(z)=1+\sum_{i\ge1}q_iz^{2i}$ one associates the
multiplicative genus $\varphi_Q$: for a closed oriented $4k$--manifold $M$ with formal
Pontryagin roots $y_1,\dots$ (so $p(TM)=\prod(1+y_j)$, $y_j=\text{(Chern root)}^2$ for
$\mathbb{CP}^n$), $\varphi_Q(M)=\big\langle\prod_jQ(\sqrt{y_j}),[M]\big\rangle$ is the
degree-$4k$ part of $\prod_jQ(\sqrt{y_j})$ paired with $[M]$. For $\mathbb{CP}^{2m}$,
where the $2m{+}1$ roots all equal $h^2$, this is $[h^{2m}]\,Q(h)^{2m+1}$.

\begin{theorem}[Characterization of the $L$--genus]\label{thm:char}
Among all even power series $Q$ with $Q(0)=1$ there is exactly one whose multiplicative
genus satisfies
\[
\varphi_Q(\mathbb{CP}^{2m})=[h^{2m}]\,Q(h)^{2m+1}=1
\qquad\text{for every }m\ge1,
\tag{12}\label{eq:normCP}
\]
namely $Q(z)=z\coth z$. Equivalently, in the Lagrange/residue formulation, $Q(z)=z/f^{-1}(z)$
where the compositional inverse is $f^{-1}=\tanh$, i.e. the genus whose logarithm is
$g(u)=\operatorname{arctanh}(u)=\sum_{m\ge0}\frac{u^{2m+1}}{2m+1}$.
\end{theorem}

\begin{proof}
\emph{Uniqueness and existence by triangularity.} Write $Q(z)=1+\sum_{i\ge1}q_iz^{2i}$.
For each $m\ge1$ the coefficient $[h^{2m}]Q(h)^{2m+1}$ is a polynomial in $q_1,\dots,q_m$
that is \emph{affine-linear and triangular in $q_m$}: expanding
$Q(h)^{2m+1}=\big(1+q_1h^2+\cdots\big)^{2m+1}$, the only way to produce $h^{2m}$ using the
single factor $q_mh^{2m}$ is $\binom{2m+1}{1}q_m\cdot 1^{2m}=(2m+1)q_m$, while every other
contribution to $[h^{2m}]$ uses only $q_1,\dots,q_{m-1}$. Hence
\[
[h^{2m}]Q(h)^{2m+1}=(2m+1)\,q_m+P_m(q_1,\dots,q_{m-1})
\]
for a universal polynomial $P_m$. The system \eqref{eq:normCP}, namely
$(2m+1)q_m+P_m=1$ for $m=1,2,3,\dots$, therefore determines $q_1,q_2,\dots$ uniquely and
recursively: $q_m=\big(1-P_m(q_1,\dots,q_{m-1})\big)/(2m+1)$. So at most one such $Q$
exists, and it exists.

\emph{Identification.} It remains to check that $Q(z)=z\coth z$ solves \eqref{eq:normCP}.
This is the residue computation already carried out in Proposition~\ref{prop:values}:
with $u=\tanh h$ (a local biholomorphism near $0$ with $du=(1-u^2)\,dh$),
\[
[h^{2m}]\Big(\frac{h}{\tanh h}\Big)^{2m+1}
=\frac{1}{2\pi i}\oint\frac{dh}{(\tanh h)^{2m+1}}
=\frac{1}{2\pi i}\oint\frac{du}{u^{2m+1}(1-u^2)}
=[u^{2m}]\frac{1}{1-u^2}=1 .
\]
Thus $z\coth z$ is the unique solution. The logarithm statement is the Hirzebruch
characteristic-series/residue dictionary: the genus with characteristic series $Q(z)=z/f^{-1}(z)$
has $\varphi_Q(\mathbb{CP}^{n})=[z^{n}]\,(f^{-1}(z)/z)^{-(n+1)}$ and equals $1$ on all even
$\mathbb{CP}^{2m}$ iff $f^{-1}(z)=\tanh z$; here $g=(f^{-1})^{-1}=\operatorname{arctanh}$
with the displayed series.
\end{proof}

\noindent\emph{Geometric meaning of \eqref{thm:char}.} The complex projective spaces are
the building blocks of positive signature ($\tau(\mathbb{CP}^{2m})=1$ by
Proposition~\ref{prop:sig}), and they rationally generate all Pontryagin numbers
(Theorem~\ref{thm:span}). The normalisation \eqref{eq:normCP} says exactly that the genus
$\varphi_Q$ is calibrated to agree with the signature on every building block; the
functional equation of $\coth$ is the analytic identity that makes this calibration
internally consistent (Theorem~\ref{thm:funeq}). The residue normalisation
$\operatorname{Res}_{z=0}\coth z=1$ is the infinitesimal form: a single transverse
intersection point contributes $+1$ to the middle intersection form.

\subsection*{Novikov additivity (proved), the geometric face of oddness}

The oddness and simple-pole structure of $\coth$ are the analytic shadow of additivity of
the signature under gluing. We now \emph{prove} the relevant additivity by the same
Poincar\'e--Lefschetz/Witt-index machinery used in Theorem~\ref{thm:cobinv}, computing the
correction term rather than waving at it.

Let $W$ be a compact oriented $(4k{+}1)$--manifold with boundary, and let $N^{4k}=\partial W$.
Define the \emph{signature of $W$}, $\tau(W)$, to be the signature of the symmetric form
on the image $I_W:=\operatorname{im}\big(H^{2k}(W,N)\to H^{2k}(W)\big)$ given by cup product
paired with $[W,N]$; this is well defined and the form is nondegenerate on $I_W$ by
Poincar\'e--Lefschetz duality (the radical of the cup form on $H^{2k}(W)$ is exactly the
kernel of $H^{2k}(W)\to H^{2k}(N)$, and $I_W$ is a complement to that radical's role; cf.
the half-lives/half-dies analysis of Theorem~\ref{thm:cobinv}).

\begin{theorem}[Novikov additivity for a separating hypersurface]\label{thm:novikov}
Let $M$ be a closed oriented $4k$--manifold and $N^{4k-1}\subset M$ a separating closed
hypersurface decomposing $M=M_+\cup_N M_-$, where $M_\pm$ are compact oriented
$4k$--manifolds with $\partial M_+=N=-\partial M_-$ (opposite induced orientations). Then
\[
\tau(M)=\tau(M_+)+\tau(M_-).
\tag{13}\label{eq:novikov}
\]
\end{theorem}

\begin{proof}
All cohomology is real. We use two standard linear-algebra facts about a nondegenerate
symmetric form $\langle\,,\rangle$ on a finite-dimensional real space $V$.

\smallskip\noindent\textbf{(W1) Witt reduction by an isotropic subspace.} If
$R\subseteq V$ is totally isotropic, then $R\subseteq R^\perp$, the form descends to a
nondegenerate symmetric form on $R^\perp/R$, and
$\operatorname{sig}\langle\,,\rangle_V=\operatorname{sig}\big(R^\perp/R\big)$. (Choose a
complement: $V\cong (R\oplus R')\perp U$ with $R\oplus R'$ a sum of hyperbolic planes of
signature $0$ and $U\cong R^\perp/R$; signatures add.)

\smallskip\noindent\textbf{(W2) Orthogonal sum.} The signature of an orthogonal direct
sum is the sum of the signatures.

\smallskip
Now consider the Mayer--Vietoris sequence of the closed cover $M=M_+\cup M_-$, with
$M_+\cap M_-\simeq N$ (a collar identification), in degree $2k$:
\[
H^{2k-1}(N)\xrightarrow{\ \partial\ }H^{2k}(M)
\xrightarrow{\ \rho:=(i_+^\ast,i_-^\ast)\ }H^{2k}(M_+)\oplus H^{2k}(M_-)
\xrightarrow{\ r:=j_+^\ast-j_-^\ast\ }H^{2k}(N).
\]
Let $V=H^{2k}(M)$ with intersection form $\langle\,,\rangle_M$, nondegenerate of signature
$\tau(M)$. Put $R:=\operatorname{im}\partial=\ker\rho$.

\emph{Step 1: $R$ is totally isotropic.} Let $x,y\in R$. By exactness $x,y$ restrict to
$0$ in $H^{2k}(M_\pm)$, so each lifts to a relative class; in particular $x$ is in the
image of $H^{2k}(M,M_-)\to H^{2k}(M)$ and $y$ in the image of $H^{2k}(M,M_+)\to H^{2k}(M)$
(the two pieces of the connecting map). Then $x\smile y$ lies in the image of
$H^{4k}(M,M_-\cup M_+)=H^{4k}(M,M)=0$, so $x\smile y=0$ in $H^{4k}(M)$ and
$\langle x,y\rangle_M=0$. Hence $R$ is totally isotropic. By (W1),
$\tau(M)=\operatorname{sig}(R^\perp/R)$ and $R^\perp/R\cong \operatorname{im}\rho$ as a
quadratic space (since $R=\ker\rho$ and the descended form on $V/R$ restricts to the
nondegenerate form on $R^\perp/R$).

\emph{Step 2: $\operatorname{im}\rho$ is the orthogonal sum of the two boundary forms.}
By exactness, $\operatorname{im}\rho=\ker r=\{(a_+,a_-):j_+^\ast a_+=j_-^\ast a_-\ \text{in}\ H^{2k}(N)\}$.
Consider the two subspaces
\[
J_+:=\operatorname{im}\big(H^{2k}(M_+,N)\to H^{2k}(M_+)\big)\hookrightarrow\operatorname{im}\rho,
\qquad
J_-:=\operatorname{im}\big(H^{2k}(M_-,N)\to H^{2k}(M_-)\big)\hookrightarrow\operatorname{im}\rho,
\]
embedded as $(a_+,0)$ and $(0,a_-)$ respectively (these land in $\ker r$ since classes
relative to $N$ restrict to $0$ on $N$). On $\operatorname{im}\rho$ the form induced from
$\langle\,,\rangle_M$ is computed by pulling back to $M_\pm$ and pairing with the relative
fundamental classes $[M_\pm,N]$: for $(a_+,a_-),(b_+,b_-)\in\ker r$,
\[
\langle a,b\rangle_M=\langle a_+\smile b_+,[M_+,N]\rangle+\langle a_-\smile b_-,[M_-,N]\rangle,
\]
because $[M]=[M_+,N]\cup_N[M_-,N]$ and the cross terms vanish (a class supported relative
to $N$ on $M_+$ cup a class on $M_-$ lies in $H^{4k}(M,M)=0$). Thus the form on
$\operatorname{im}\rho$ is the orthogonal sum of the cup forms on the $M_+$-part and the
$M_-$-part; the nondegenerate quotient $R^\perp/R$ identifies orthogonally with
$I_{M_+}\perp I_{M_-}$, the nondegenerate boundary forms of Definition (signature
$\tau(M_\pm)$). By (W2),
\[
\tau(M)=\operatorname{sig}(R^\perp/R)=\operatorname{sig}(I_{M_+})+\operatorname{sig}(I_{M_-})
=\tau(M_+)+\tau(M_-).
\]
The correction obstructing naive additivity is exactly $R=\operatorname{im}\partial$, the
totally isotropic image of the connecting map, which is signature-neutral by (W1): this is
the geometric content of the residue $1$/oddness of $\coth$ --- contributions concatenate
and the gluing locus contributes $0$.
\end{proof}

\begin{remark}
This is the \emph{separating} (Novikov) case, where the form on $H^{2k}(N)$ plays no
obstructing role; in non-separating gluings Wall's non-additivity term appears, governed by
a Maslov triple index. The separating case suffices for the geometric interpretation here
and is exactly the case mirrored by the simple additive pole structure of $\coth$.
\end{remark}

\subsection*{Multiplicativity in fibre bundles (stated and cited)}

The product case is elementary and proved:

\begin{proposition}[Product multiplicativity]\label{prop:prodmult}
For closed oriented $F,B$ of dimensions divisible by $4$,
$\tau(F\times B)=\tau(F)\tau(B)$ and $L[F\times B]=L[F]\,L[B]$.
\end{proposition}

\begin{proof}
This is fact (I): by K\"unneth the intersection form of $F\times B$ is the tensor product
of the two forms, and the signature of a tensor product of nondegenerate symmetric real
forms is the product of the signatures (diagonalise both; eigenvalue signs multiply). For
$L$, multiplicativity of multiplicative sequences and $p(T(F\times B))=p(TF)\times p(TB)$
give $L[F\times B]=L[F]\,L[B]$.
\end{proof}

The deepest geometric incarnation of the functional equation is the
\emph{Chern--Hirzebruch--Serre} multiplicativity of the signature in fibre bundles. We
state it precisely with hypotheses, and \emph{cite} it: its proof uses the Leray--Serre
spectral sequence structure of the fibration and is not elementary in the present sense, so
we do not claim a proof.

\begin{theorem}[Chern--Hirzebruch--Serre, cited]\label{thm:CHS}
Let $F\hookrightarrow E\xrightarrow{\pi}B$ be a smooth oriented fibre bundle of closed
oriented manifolds with $\dim E\equiv0\ (\mathrm{mod}\ 4)$, such that the fundamental group
$\pi_1(B)$ acts trivially on $H^\ast(F;\mathbb Q)$ (e.g. $B$ simply connected). Then
$\tau(E)=\tau(F)\,\tau(B)$.
\end{theorem}

The link to \eqref{eq:funeq} is the \emph{formal group law}: a multiplicative genus
$\varphi_Q$ has an associated formal group; for $Q(z)=z\coth z$ the group law is the
multiplicative one $F(x,y)=\dfrac{x+y}{1+xy}$ (the addition law of $\tanh$, equivalently of
$\coth$ via \eqref{eq:funeq}), and it is exactly this group structure that makes the genus
behave multiplicatively under the fibration operation. Among normalised even genera, the
addition law \eqref{eq:funeq} singles out $z\coth z$ as the genus compatible with this
fibrewise composition; the elementary special case ($B,F$ a product) is
Proposition~\ref{prop:prodmult}.

\subsection*{Summary of the geometric interpretation}

\begin{itemize}
\item[(a)] \textbf{Residue normalisation} $\operatorname{Res}_{z=0}\coth z=1$
(Theorem~\ref{thm:funeq}) $\Leftrightarrow$ a transverse positive intersection point
contributes $+1$; it calibrates $\varphi_Q$ to equal $\tau$ on the building blocks
$\mathbb{CP}^{2m}$ (Theorem~\ref{thm:char}, Propositions~\ref{prop:sig},
\ref{prop:values}).
\item[(b)] \textbf{Oddness + simple-pole structure} $\Leftrightarrow$ \textbf{Novikov
additivity} \eqref{eq:novikov}, \emph{proved} in Theorem~\ref{thm:novikov}: signatures of
glued pieces add and the gluing locus is a totally isotropic, signature-neutral
correction.
\item[(c)] \textbf{Addition law} $\coth(z+w)=\frac{1+\coth z\coth w}{\coth z+\coth w}$
(Theorem~\ref{thm:funeq}) $\Leftrightarrow$ the multiplicative formal group law underlying
\textbf{multiplicativity in fibre bundles} (Theorem~\ref{thm:CHS}, cited), whose elementary
shadow is product multiplicativity (Proposition~\ref{prop:prodmult}). The
characterization Theorem~\ref{thm:char} makes precise the sense in which $z\coth z$ --- and
no other normalised genus --- is the one selected by these requirements.
\end{itemize}
This completes the geometric interpretation requested in part (b): the analytic functional
equation (Theorem~\ref{thm:funeq}) and the characterization (Theorem~\ref{thm:char}) are
proved elementarily; the additivity incarnation (Theorem~\ref{thm:novikov}) is proved in
full; and the fibre-bundle incarnation (Theorem~\ref{thm:CHS}) is stated with verified
hypotheses and cited, with its elementary special case proved.


\end{document}
